% Generated draft: requires mathematical and source review.
\documentclass[11pt]{article}
\usepackage[margin=25mm]{geometry}
\usepackage{fontspec}
\setmainfont{Latin Modern Roman}
\usepackage{amsmath,amssymb,mathtools,unicode-math}
\setmathfont{Latin Modern Math}
\usepackage{enumitem,needspace,xurl,hyperref}
\setlength{\parindent}{0pt}
\setlength{\parskip}{5pt}
\setlength{\emergencystretch}{3em}
\newcommand{\N}{\mathbb{N}}
\newcommand{\Z}{\mathbb{Z}}
\newcommand{\Q}{\mathbb{Q}}
\newcommand{\R}{\mathbb{R}}
\newcommand{\C}{\mathbb{C}}
\newcommand{\F}{\mathbb{F}}
\newcommand{\A}{\mathbb{A}}
\newcommand{\PP}{\mathbb{P}}
\newcommand{\E}{\mathbb{E}}
\newcommand{\eps}{\varepsilon}
\newcommand{\dd}{\,\mathrm{d}}
\begin{document}
\section*{UnsolvedMath ID: 30005403}
\textbf{Source problem code:} catalogue entry
\subsection{Definitions and mathematical statement}
Let \(W_{\Q_p}\) denote the Weil group of \(\Q_p\). The supplied excerpt considers a group \(G\), a torus \(T\), a toral parameter \(\phi_T\), and a parameter \(\phi\) induced from \(\phi_T\). It does not specify the full hypotheses on these objects, the coefficient field, or the precise parameter maps. These data must therefore remain unresolved rather than being assigned additional hypotheses.

For each coroot \(\alpha\) occurring in the source's convention, write \(\chi_\alpha=\alpha\circ\phi_T\) for the resulting character of \(W_{\Q_p}\). The source defines \(\phi_T\) to be \emph{generic} when
\[
\forall\alpha,\qquad \chi_\alpha\not\cong\mathbf{1}\quad\text{and}\quad\chi_\alpha\not\cong|\cdot|,
\]
where \(\mathbf{1}\) is the trivial character and \(|\cdot|\) is the norm character in the source's normalization. The source also states the equivalent condition
\[
\forall\alpha,\qquad R\Gamma(W_{\Q_p},\chi_\alpha)\simeq 0.
\]
The relevant coroot convention, coefficient category, and norm normalization are not specified in the excerpt; the asserted equivalence is retained as a source assertion subject to those missing conventions.

\textbf{Intended assertion, recoverable only schematically.} For a parameter \(\phi\) induced from a generic toral parameter \(\phi_T\), the proposed sheaf \(S_\phi\) should be a Hecke eigensheaf with eigenvalue \(\phi\). Here the eigensheaf condition requires compatible Hecke eigenisomorphisms, but neither the relevant Hecke functors nor the construction of \(S_\phi\) is supplied. Thus this sentence does not specify a complete mathematical conjecture.

The excerpt explicitly calls the unrestricted assertion naive: representations of non-quasi-split groups can have semisimplified parameters factoring through a maximal torus and are expected to contribute to eigensheaf stalks. It gives the trivial representation of the units of a quaternion division algebra as an example, although the extracted algebra notation is damaged. It further asserts that genericity excludes \(\phi\) from being the semisimplification of a parameter with nontrivial monodromy. Its restriction on possible eigensheaf stalks additionally assumes expected properties of the Fargues--Scholze correspondence and refers to an earlier stalk description absent from the excerpt.

\textbf{Conditional result reported by the source.} The excerpt attributes to its reference [6, Theorem 1.17] the eigensheaf assertion for generic toral parameters, assuming certain properties of the Fargues--Scholze local Langlands correspondence and possibly additional constraints on \(\phi_T\). Those properties and constraints are not enumerated here, so a precise conditional theorem cannot be reconstructed.

The discussion uses the moduli space \(\operatorname{Bun}_B\) of parabolic \(B\)-structures on \(G\)-bundles, a Drinfeld compactification, and Eisenstein constructions applied to the toral eigensheaf \(S_{\phi_T}\). It distinguishes pushforward from the open space from pushforward from its compactification, the latter requiring an intersection-cohomology kernel. The extraction loses the distinguishing notation and ends mid-sentence. Consequently it does not determine which construction defines \(S_\phi\), or a standalone unresolved assertion about it. The registry explicitly marks the statement as unrecoverable; this draft preserves that limitation.
\subsection{Short English statement}
The intended claim is that a sheaf associated with a generic toral parameter is a Hecke eigensheaf for the induced parameter. The excerpt rejects an unrestricted naive formulation and reports a conditional theorem, but omits the sheaf construction and essential hypotheses. It therefore does not support a precise standalone open problem without recovery and mathematical review of the source.
\subsection{Sources}
\begin{itemize}
\item Local UnsolvedMath registry, permanent ID \(30005403\), source code \texttt{catalogue entry}, title ``Naive Local--Global Compatibility Formulation.'' The supplied record has proposed year 2023, no named proposer, status \texttt{partially\_solved}, and statement status \texttt{unrecoverable}. Its supplied verification note says that the official report uses the naive formulation as exposition before a conditional theorem and that the extraction does not preserve a standalone unresolved claim.
\item \emph{Arithmetic of Shimura Varieties} (2023), Oberwolfach Reports, pp. 261--326, DOI 10.4171/owr/2023/5. Supplied workshop-report provenance: \url{https://doi.org/10.4171/owr/2023/5}.
\item Supplied literature reference: ``Hamann, geometric Eisenstein series,'' \url{https://arxiv.org/abs/2209.08175}. The excerpt cites [6, Theorem 1.17], but does not reproduce the numbered bibliography needed to establish that reference's identity from the excerpt alone.
\item Supplied literature reference: ``OWR statement and conditional theorem,'' \url{https://ems.press/content/serial-article-files/46998?nt=1}.
\item The excerpt attributes classical geometric Langlands background to Braverman--Gaitsgory--Laumon, references [2, 7], without supplying their bibliographic entries.
\item Supplied literature assessment, checked \textbf{2026-08-22}: the generic-toral eigensheaf claim is described as proved in restricted settings or under stated Fargues--Scholze compatibility assumptions, with the unconditional general assertion described as remaining unproved. This is supplied dated curation derived from the OWR record, not an independently verified mathematical or current-literature assessment.
\end{itemize}
\end{document}
